\section{Estudio de Factibilidad y Gestión de Riesgos}

El presente estudio analiza la viabilidad del proyecto desde cuatro dimensiones
complementarias: técnica, operativa, económica y legal. El objetivo es demostrar,
con evidencia cuantitativa y referencias documentadas, que el proyecto puede
construirse con los recursos disponibles, ser adoptado por sus usuarios finales,
sostenerse económicamente y operar dentro del marco normativo vigente en México.

La metodología aplicada combina: (1) revisión de documentación oficial de las
plataformas tecnológicas seleccionadas; (2) análisis de estadísticas de adopción
digital publicadas por el INEGI y organismos de la industria; (3) modelado de
costos basado en tarifas públicas vigentes; y (4) revisión del marco normativo
mexicano aplicable al manejo de datos personales, transacciones electrónicas y
software \cite{nextjs, supabase_pricing, prisma_docs, state_of_js_2023, stackoverflow_2023_results}.

% -----------------------------------------------------------
\FloatBarrier
\subsection{Factibilidad Técnica}
% -----------------------------------------------------------

La factibilidad técnica evalúa la disponibilidad y madurez de las herramientas,
lenguajes, frameworks e infraestructura necesarios para el desarrollo e implantación
del sistema.

\subsubsection{Ecosistema de Desarrollo: Madurez y Adopción}

La selección del stack tecnológico se fundamenta en el criterio de madurez
comprobada, adopción masiva en la industria y soporte activo de la comunidad.
El uso de tecnologías ampliamente adoptadas reduce el riesgo técnico, amplía el
pool de desarrolladores disponibles y garantiza continuidad del soporte a largo
plazo. La Tabla~\ref{tab:stack} resume las tecnologías principales con sus métricas
de adopción.

\begin{table}[H]
\caption{Stack Tecnológico: Madurez y Adopción}
\label{tab:stack}
\centering
\small
\setlength{\tabcolsep}{3pt}
\begin{tabular}{>{\raggedright\arraybackslash}p{2.3cm}
                >{\raggedright\arraybackslash}p{3cm}
                >{\raggedright\arraybackslash}p{1.5cm}
                >{\raggedright\arraybackslash}p{6cm}
                >{\raggedright\arraybackslash}p{2.5cm}}
\hline
\thead{Tecnología} & \thead{Categoría} & \thead{Versión} & \thead{Adopción} & \thead{Licencia} \\
\hline
Next.js 14 & Framework frontend + API & 14.x & 44\% de desarrolladores lo usan activamente \cite{state_of_js_2023} & MIT \\ \rowcolor{altrow}
React 18 & Biblioteca UI & 18.x & 84\% de los desarrolladores web la usan \cite{state_of_js_2023} & MIT \\
TypeScript 5 & JS tipado & 5.x & 38.87\% de desarrolladores la usan \cite{stackoverflow_2023_results} & Apache 2.0 \\ \rowcolor{altrow}
Supabase & BaaS: PostgreSQL + Auth + Realtime & GA 2023 & +65,000 estrellas en GitHub; 200,000+ proyectos activos \cite{supabase_pricing} & Apache 2.0 \\
PostgreSQL 15 & SGBD relacional & 15.x & 45.55\% de popularidad \cite{stackoverflow_2023_results} & PostgreSQL License \\ \rowcolor{altrow}
Prisma ORM 5 & ORM para Node.js/TS & 5.x & ORM preferido en ecosistema Next.js/TS \cite{prisma_docs} & Apache 2.0 \\
Vercel & Hosting serverless & PaaS & Plataforma l\'ider para despliegue Next.js \cite{vercel_pricing, nextjs} & Propietario \\ \rowcolor{altrow}
Node.js 20 LTS & Runtime JS servidor & 20.x LTS & framework m\'as usado (42.65\%) \cite{stackoverflow_2023_results} & MIT \\
\hline
\end{tabular}
\end{table}

\subsubsection{Análisis de Componentes Críticos}

\noindent\textbf{Vercel como plataforma de hosting.}
Vercel es la plataforma de referencia para despliegue de aplicaciones Next.js,
siendo el creador y principal mantenedor del framework. Ofrece escalado automático,
red CDN global en más de 100 edge locations, deploy automático por rama de Git y
ambientes de Preview y Producción nativos. El plan Pro (\$20 USD/mes) elimina las
restricciones del plan gratuito, incluye 60 segundos de tiempo de ejecución para
funciones serverless y soporte para dominio personalizado con SSL gestionado
automáticamente \cite{vercel_pricing, nextjs}.

\noindent\textbf{Supabase como backend-as-a-service.}
Supabase provee una instancia PostgreSQL administrada junto con servicios de
autenticación, Realtime (basado en LISTEN/NOTIFY de PostgreSQL), almacenamiento
de archivos y una API REST auto-generada. Su elección está fundamentada en que
PostgreSQL es el sistema de base de datos m\'as popular (45.55\% de adopci\'on en 2023) \cite{stackoverflow_2023_tech, stackoverflow_2023_results}.
El plan Pro (\$25 USD/mes) incluye 8~GB de almacenamiento, backups diarios
automatizados y acceso a conexiones ilimitadas a través de PgBouncer \cite{supabase_pricing}.

\noindent\textbf{Prisma ORM como capa de acceso a datos.}
Prisma provee un cliente de base de datos tipado generado automáticamente a partir
del schema definido en su DSL. Esto garantiza seguridad de tipos en tiempo de
compilación, previene SQL Injection por diseño al utilizar consultas parametrizadas
y facilita las migraciones de schema con control de versiones. En entornos
serverless, Prisma Accelerate resuelve el problema de agotamiento de conexiones
al actuar como proxy entre las funciones de Vercel y la instancia PostgreSQL
de Supabase \cite{prisma_docs, prisma_accelerate}.

La Tabla~\ref{tab:hw_usuarios} detalla los requerimientos mínimos de hardware
para cada tipo de usuario final.

\begin{table}[H]
\caption{Requerimientos de Hardware por Tipo de Usuario}
\label{tab:hw_usuarios}
\centering
\small
\setlength{\tabcolsep}{3pt}
\begin{tabular}{>{\raggedright\arraybackslash}p{2.3cm}
                >{\raggedright\arraybackslash}p{2.5cm}
                >{\raggedright\arraybackslash}p{5.5cm}
                >{\raggedright\arraybackslash}p{5cm}}
\hline
\thead{Tipo de usuario} & \thead{Dispositivo} & \thead{Requerimiento mínimo} & \thead{Observación} \\
\hline
Comensal / Anfitrión & Smartphone personal & Chrome $\geq$90, Firefox $\geq$88 o Safari $\geq$14. Cámara para QR. & El 87.3\% de mexicanos de 18--44 años usa smartphone con internet \cite{inegi_endutih_2022}. \\ \rowcolor{altrow}
Mesero & Tablet del restaurante & Chrome $\geq$90, Firefox $\geq$88 o Safari $\geq$14. Pantalla $\geq$8 pulgadas recomendada. & Tablet de gama media ($\sim{\$4,500 MXN}$). \\
Chef & Tablet del restaurante & Pantalla $\geq$8 pulgadas recomendada para tablero de órdenes. & Tablet o monitor con PC básica. \\ \rowcolor{altrow}
Gerente / Admin & PC, laptop o tablet & Navegador moderno. Pantalla $\geq$11 pulgadas recomendada para Dashboard. & Equipo existente en la mayoría de negocios. \\
Servidor (VPS) & VPS en la nube & 2 vCPU, 4~GB RAM, 40~GB SSD. Ubuntu 22.04~LTS. & Hetzner CX22: \$6.37~USD/mes. \\
\hline
\end{tabular}
\end{table}

\subsubsection{Disponibilidad de Conocimientos Técnicos}

El stack seleccionado (TypeScript + Next.js + PostgreSQL + React) es el más
demandado en el mercado de desarrollo web en México y Latinoamérica. TypeScript
ha sido el lenguaje más deseado por desarrolladores tres años consecutivos
(2021--2023) \cite{state_of_js_2023}; Next.js ocupa el primer lugar entre frameworks full-stack
deseados \cite{lftr}; y PostgreSQL es la base de datos relacional más popular globalmente \cite{state_of_js_2023}.
Plataformas como LinkedIn, OCC Mundial y Computrabajo registran cientos de
ofertas activas con este perfil mensualmente.

\subsubsection{Restricciones Técnicas y Mitigación}

La Tabla~\ref{tab:restricciones_tec} resume las restricciones técnicas identificadas
y las soluciones disponibles para cada una.

\begin{table}[H]
\caption{Restricciones Técnicas Identificadas y Mitigación}
\label{tab:restricciones_tec}
\centering
\small
\setlength{\tabcolsep}{3pt}
\begin{tabular}{>{\raggedright\arraybackslash}p{4cm}
                >{\raggedright\arraybackslash}p{4cm}
                >{\raggedright\arraybackslash}p{7.5cm}}
\hline
\thead{Restricción técnica} & \thead{Impacto potencial} & \thead{Solución técnica disponible} \\
\hline
Límite de conexiones PostgreSQL en funciones serverless &
  Pool de conexiones agotado en horas pico &
  Prisma Accelerate o PgBouncer (integrado en Supabase Pro). Solución probada y documentada \cite{prisma_docs, prisma_accelerate}. \\ \rowcolor{altrow}
Timeout de funciones Vercel (60~s en plan Pro) &
  Operaciones largas de pago podrían fallar &
  Diseño de endpoints con respuesta $<$5~s. Procesamiento asíncrono para operaciones largas. \\
Latencia de Supabase Realtime bajo carga alta &
  Retrasos en actualización del tablero del Chef &
  Supabase Realtime plan Pro soporta hasta 500 conexiones simultáneas. Escalable bajo demanda \cite{supabase_pricing}. \\ \rowcolor{altrow}
Acceso a internet requerido en todo momento &
  Sin conectividad, la app no funciona &
  Restricción asumida y documentada. El contexto de uso (restaurante) garantiza conectividad WiFi. \\
\hline
\end{tabular}
\end{table}

\noindent\textbf{Veredicto:} El stack tecnológico seleccionado es maduro,
ampliamente adoptado en la industria y con soporte activo. Todas las restricciones
técnicas identificadas tienen soluciones documentadas y probadas en producción
por proyectos de escala comparable. La plataforma puede construirse con el
conocimiento técnico disponible en el mercado mexicano. \textbf{FACTIBLE}.

% -----------------------------------------------------------
\FloatBarrier
\subsection{Factibilidad Operativa}
% -----------------------------------------------------------

La factibilidad operativa evalúa la capacidad de los usuarios finales para
adoptar y operar el sistema propuesto, considerando el perfil tecnológico de
la población objetivo en México.

\subsubsection{Perfil Tecnológico de los Usuarios Finales}

\noindent\textbf{Comensales: adopción de smartphones y QR en México.}
El usuario comensal interactúa con el sistema únicamente a través del navegador
de su dispositivo móvil. Según la Encuesta Nacional sobre Disponibilidad y Uso de
Tecnologías de la Información en los Hogares (ENDUTIH) 2022 del INEGI, el 88.6\%
de la población mexicana de 25 a 34 años es usuaria de internet, y el 89.3\%
accede a internet principalmente desde un teléfono celular \cite{inegi_endutih_2022}. En el segmento de
35 a 44 años, la penetración de internet es del 78.9\%.

El uso de c\'odigos QR para acceder a men\'us digitales se consolid\'o durante
y despu\'es de la pandemia de COVID-19. De acuerdo con la National Restaurant
Association (2023), el 73\% de los comensales de la Generaci\'on~Z y el 46\%
de los Baby~Boomers muestran inter\'es en usar c\'odigos QR para men\'us de
restaurante \cite{nra_state_2023}. A nivel nacional, la alta penetraci\'on de smartphones en
M\'exico \cite{inegi_endutih_2022} respalda la viabilidad de este canal como punto de entrada al sistema.
Esto implica que la curva de aprendizaje para el flujo de activaci\'on de mesa
virtual es m\'inima.

\noindent\textbf{Personal operativo: capacitación requerida.}
Las interacciones del sistema han sido diseñadas para minimizar la curva de
aprendizaje: el mesero solo necesita inicializar mesas y entregar un QR; el chef
solo visualiza un tablero de órdenes; el gerente administra un panel web estándar.
La Tabla~\ref{tab:capacitacion} detalla la curva de aprendizaje estimada por rol.

\begin{table}[H]
\caption{Capacitación Requerida por Rol}
\label{tab:capacitacion}
\centering
\small
\setlength{\tabcolsep}{3pt}
\begin{tabular}{>{\raggedright\arraybackslash}p{2.5cm}
                >{\raggedright\arraybackslash}p{4.5cm}
                >{\raggedright\arraybackslash}p{2.5cm}
                >{\raggedright\arraybackslash}p{6.5cm}}
\hline
\thead{Rol} & \thead{Flujos principales} & \thead{Curva de aprendizaje} & \thead{Capacitación requerida} \\
\hline
Mesero & Inicializar mesa, entregar QR, visualizar estatus, liberar mesa. & 2--4 horas & Sesión práctica de 1~hora + manual de 1~página. \\ \rowcolor{altrow}
Chef & Visualizar tablero de órdenes, actualizar estatus de orden. & 1--2 horas & Sesión práctica de 30~minutos. Interfaz de una sola pantalla. \\
Gerente de Sucursal & Todo lo anterior + gestión de menú, usuarios y dashboard. & 4--8 horas & Capacitación de 2~horas + documentación de usuario. \\ \rowcolor{altrow}
Gerente de Zona / Admin & Gestión multi-sucursal, zonas, analytics. & 6--10 horas & Capacitación de 3~horas + acceso a modo sandbox. \\
Comensal / Anfitrión & Escanear QR, crear contraseña, ordenar, pagar. & 0~horas & Interfaz autoexplicativa con flujo guiado y apoyo del mesero. \\
\hline
\end{tabular}
\end{table}

\subsubsection{Análisis de Cambio Organizacional}

La Tabla~\ref{tab:cambio_org} evalúa los factores de cambio que implica la
implantación del sistema y los mecanismos de mitigación propuestos.

\begin{table}[H]
\caption{Análisis de Cambio Organizacional}
\label{tab:cambio_org}
\centering
\small
\setlength{\tabcolsep}{3pt}
\begin{tabular}{>{\raggedright\arraybackslash}p{3.5cm}
                >{\raggedright\arraybackslash}p{5.5cm}
                >{\raggedright\arraybackslash}p{7cm}}
\hline
\thead{Factor de cambio} & \thead{Impacto para el personal} & \thead{Mitigación propuesta} \\
\hline
El comensal genera su propia orden &
  El mesero deja de tomar órdenes en papel. Puede percibirse como reducción de su rol. &
  Comunicar que el mesero se libera para atención de mayor valor. No implica reducción de plantilla. \\ \rowcolor{altrow}
Nuevo flujo de activación de mesa &
  Paso adicional al proceso actual (inicializar mesa virtual antes de sentar comensales). &
  El flujo toma menos de 60~segundos. Onboarding interactivo integrado en la app. \\
Chef visualiza tablero digital &
  Cambio de medio: papel $\rightarrow$ pantalla. Puede generar resistencia en personal de mayor edad. &
  Interfaz con fuente grande. Notificación sonora de nuevas órdenes. Período de transición paralela recomendado. \\ \rowcolor{altrow}
Comensales pagan con tarjeta desde su teléfono &
  Comensales que prefieren efectivo pueden percibir fricción. &
  El sistema no bloquea el pago en caja como alternativa. El mesero puede registrar cierres manuales. \\
\hline
\end{tabular}
\end{table}

\subsubsection{Contexto de la Industria Restaurantera en México}

Según la Cámara Nacional de la Industria de Restaurantes y Alimentos Condimentados
(CANIRAC), el sector restaurantero en M\'exico concentra m\'as de 700{,}000 unidades
econ\'omicas seg\'un datos del DENUE del INEGI \cite{inegi_denue_2024}. El sector ha experimentado una aceleraci\'on en
la adopción de tecnología post-pandemia: el 68\% de los restaurantes de servicio
completo en zonas urbanas ya utilizan algún sistema digital de gestión de pedidos
o punto de venta \cite{gonzalez_canirac_2024}.

\noindent\textbf{Veredicto:} El perfil tecnológico de los usuarios finales es
compatible con los requerimientos del sistema. El 89.3\% de los mexicanos accede
a internet desde smartphone \cite{inegi_endutih_2022}, y el 64\% ha usado un QR para acceder a un menú
\cite{nra_state_2023}. Los flujos del sistema están diseñados para minimizar la curva
de aprendizaje. La resistencia al cambio en el personal operativo es manejable
con capacitación mínima y transición gradual. \textbf{FACTIBLE}.

% -----------------------------------------------------------
\FloatBarrier
\subsection{Factibilidad Económica}

\subsubsection{Metodología de Estimación de Costos}

El presente estudio de factibilidad econ\'omica aplica la metodolog\'ia de
\textbf{Estimaci\'on Bottom-Up}, definida por el Project Management Institute
como la descomposici\'on del proyecto en actividades individuales con estimaci\'on
de costo independiente para cada una, que posteriormente se suman para obtener
el presupuesto total \cite{stps_observatorio}. Este m\'etodo es id\'oneo para el presente proyecto
porque los m\'odulos (OE-01 a OE-13) est\'an claramente delimitados y permiten
estimar horas de desarrollo por objetivo espec\'ifico con alta precisi\'on.

Los par\'ametros base utilizados son: tarifa de \$500~MXN/hora para un desarrollador
\textit{full-stack senior} con el perfil TypeScript~+~Next.js~+~Supabase, consistente
con el rango medio del mercado mexicano de desarrollo de software en 2026 seg\'un
el Observatorio Laboral de la STPS \cite{stps_observatorio}; y un tipo de cambio referencial de
\$17.50~MXN/USD conforme a Banxico. Ambos par\'ametros son editables en el archivo
de gesti\'on de costos adjunto (\textit{Gesti\'on\_Costos.xlsx}).

\subsubsection{Estimación de Costos de Desarrollo (Inversión Inicial)}

La Tabla~\ref{tab:costos_dev} resume los costos de desarrollo por m\'odulo,
derivados del desglose bottom-up de los 13 objetivos espec\'ificos y las
actividades de la Fase~1 (an\'alisis y dise\~no). El detalle completo, con
horas por actividad y f\'ormulas editables, se encuentra en la hoja
\textit{2.~Estimaci\'on Bottom-Up} del archivo Excel adjunto.

\setstretch{1.2}
\setlength{\tabcolsep}{4pt}
\begin{longtable}{>{\raggedright\arraybackslash}p{3.5cm}
                >{\raggedright\arraybackslash}p{5.5cm}
                >{\centering\arraybackslash}p{1.6cm}
                >{\centering\arraybackslash}p{2.5cm}
                >{\centering\arraybackslash}p{2cm}}
\caption{Estimación de Costos de Desarrollo por Módulo (Bottom-Up)} \label{tab:costos_dev} \\
\hline
\thead{Módulo / Fase} & \thead{Actividades principales} &
\thead{Horas est.} & \thead{Costo RRHH\,(MXN)} & \thead{Fase} \\
\hline
\endfirsthead
\caption[]{TABLA \thetable{} (Continuación)} \\
\hline
\thead{Módulo / Fase} & \thead{Actividades principales} &
\thead{Horas est.} & \thead{Costo RRHH\,(MXN)} & \thead{Fase} \\
\hline
\endhead
\hline
\endfoot
OE-01 Infraestructura & Vercel CI/CD, schema Prisma, Supabase base & 160 & \$80,000 & 2 \\ \rowcolor{altrow}
OE-02 Estructura Org. & CRUD Cadena/Zona/Sucursal + RLS & 60 & \$30,000 & 2 \\
OE-03 Auth \& Roles & Supabase Auth + RLS policies + JWT RS256 & 100 & \$50,000 & 2 \\ \rowcolor{altrow}
OE-04 Gestión Usuarios & Alta usuarios + contrase\~na temporal & 50 & \$25,000 & 2 \\
OE-05 Cat\'alogo & CRUD categor\'ias + productos + variantes & 90 & \$45,000 & 2 \\ \rowcolor{altrow}
OE-06/07 Mesas & Mesas f\'isicas + mesa virtual + QR HMAC & 150 & \$75,000 & 2 \\
OE-08/09 \'Ordenes + Chef & Canasta, \'ordenes, tablero KDS Realtime & 170 & \$85,000 & 2 \\ \rowcolor{altrow}
OE-10/11 Liquidaci\'on & Concurrencia \texttt{SELECT FOR UPDATE} + comprobante & 140 & \$70,000 & 2 \\
OE-12 Dashboard & Pipeline Spark (Medallion) + Dashboard multi-nivel & 120 & \$60,000 & 2 \\ \rowcolor{altrow}
OE-13 Seguridad & Rate limiting + RLS audit + pruebas carga & 80 & \$40,000 & 2 \\
Pruebas QA & Unitarias + integraci\'on + carga (k6) & 80 & \$40,000 & 2 \\ \rowcolor{altrow}
Fase 1 --- An\'alisis y Dise\~no (Natalia) & Marco te\'orico, flujos, UML, BD, wireframes, TT & 660 & \$330,000 & 1 \\
Gesti\'on de Proyecto (ambos) & Ceremonias Scrum, documentaci\'on, retrospectivas & 80 & \$40,000 & 1/2 \\
\hline
\rowcolor{green}
\textbf{TOTAL RRHH} & & \textbf{1,940} & \textbf{\$970,000} & --- \\
\end{longtable}
\setstretch{1.5}

\subsubsection{Costos de Infraestructura Cloud (Operación Mensual)}

Los costos de infraestructura son recurrentes y se calculan con base en las tarifas
p\'ublicas de cada servicio. La Tabla~\ref{tab:costos_infra} presenta los costos
mensuales y anuales para el primer a\~no del proyecto (11 meses operativos).

\begin{table}[H]
\caption{Costos de Infraestructura Cloud (Año 1 — 11 Meses)}
\label{tab:costos_infra}
\centering
\small
\setlength{\tabcolsep}{3pt}
\begin{tabular}{>{\raggedright\arraybackslash}p{3cm}
                >{\raggedright\arraybackslash}p{2.5cm}
                >{\centering\arraybackslash}p{1.5cm}
                >{\centering\arraybackslash}p{1.5cm}
                >{\centering\arraybackslash}p{2.5cm}
                >{\raggedright\arraybackslash}p{5cm}}
\hline
\thead{Servicio} & \thead{Plan} & \thead{USD/mes} &
\thead{MXN/mes} & \thead{MXN\,(11 meses)} & \thead{Referencia} \\
\hline
Vercel & Pro (1 seat) & \$20 & \$350 & \$3,850 & vercel.com/pricing (2026) \cite{vercel_pricing} \\ \rowcolor{altrow}
Supabase & Pro & \$25 & \$437.50 & \$4,812.50 & supabase.com/pricing (2026) \cite{supabase_pricing} \\
VPS Hetzner CX22 & 2~vCPU / 4~GB RAM & \$6.37 & \$111.50 & \$1,226.50 & hetzner.com/cloud (2026) \\ \rowcolor{altrow}
Dominio .com.mx & Anual & \$0.83 & \$14.53 & \$159.83 & NIC M\'exico / GoDaddy (2026) \\
\hline
\rowcolor{green}
\textbf{Total infraestructura} & & \textbf{\$52.20} & \textbf{\$913.53} & \textbf{\$10,048.83} & \\
\hline
\end{tabular}
\end{table}

\subsubsection{Presupuesto Total del Proyecto (Año 1)}

La Tabla~\ref{tab:presupuesto_total} presenta el presupuesto consolidado del primer
a\~no, incluyendo la reserva de contingencia del 15\% conforme al PMI PMBOK \cite{pmi_pmbok_7}.

\begin{table}[H]
\caption{Presupuesto Total del Proyecto — Año 1}
\label{tab:presupuesto_total}
\centering
\small
\setlength{\tabcolsep}{4pt}
\begin{tabular}{>{\raggedright\arraybackslash}p{7cm}
                >{\centering\arraybackslash}p{2cm}
                >{\centering\arraybackslash}p{3.5cm}
                >{\centering\arraybackslash}p{3.5cm}}
\hline
\thead{Concepto} & \thead{Tipo} & \thead{Costo MXN} & \thead{Costo USD} \\
\hline
Recursos Humanos --- Desarrollo (Diego) & Directo/Fijo & \$620,000 & \$35,429 \\ \rowcolor{altrow}
Recursos Humanos --- An\'alisis y Dise\~no (Natalia) & Directo/Fijo & \$330,000 & \$18,857 \\
Gesti\'on de Proyecto (ambos) & Directo/Fijo & \$40,000 & \$2,286 \\ \rowcolor{altrow}
Infraestructura cloud (11 meses) & Directo/Fijo & \$10,049 & \$574 \\
Capacitaci\'on del personal (onboarding) & Directo/Fijo & \$12,000 & \$686 \\ \rowcolor{altrow}
Consultor\'ia legal (LFPDPPP + T\'erminos) & Indirecto/Fijo & \$25,000 & \$1,429 \\
Herramientas de dise\~no (Figma Pro 6 meses) & Directo/Variable & \$3,150 & \$180 \\ \rowcolor{altrow}
Herramientas de documentaci\'on y PM & Indirecto/Variable & \$1,750 & \$100 \\
\hline
\rowcolor{orange}
\textbf{Subtotal Base} & & \textbf{\$1,041,949} & \textbf{\$59,540} \\ \rowcolor{orange}
Reserva de Contingencia (15\% --- PMI PMBOK) & Buffer & \$156,292 & \$8,931 \\
\hline
\rowcolor{green}
\textbf{PRESUPUESTO TOTAL A\~NO 1} & & \textbf{\$1,198,241 MXN} & \textbf{\$68,471 USD} \\
\hline
\end{tabular}
\end{table}

\subsubsection{Análisis Comparativo: Desarrollo Propio vs. Alternativas SaaS}

Una decisi\'on relevante en el estudio econ\'omico es comparar el costo de desarrollo
propio frente a soluciones comerciales disponibles en el mercado. La
Tabla~\ref{tab:build_buy} presenta la comparaci\'on para una cadena de 3~sucursales.

\begin{table}[H]
\caption{Comparativo Build vs. Buy --- 3 Sucursales (Costos Anuales MXN)}
\label{tab:build_buy}
\centering
\small
\setlength{\tabcolsep}{3pt}
\begin{tabular}{>{\raggedright\arraybackslash}p{3.5cm}
                >{\centering\arraybackslash}p{2.5cm}
                >{\centering\arraybackslash}p{3.5cm}
                >{\centering\arraybackslash}p{2cm}
                >{\raggedright\arraybackslash}p{4.5cm}}
\hline
\thead{Solución} & \thead{Costo/suc./mes\,(MXN)} & \thead{Costo anual\,3 suc.\,(MXN)} &
\thead{Propie-dad código} & \thead{Observación} \\
\hline
\textbf{El presente sistema} & \textbf{\$304} & \textbf{\$10,049} & \textbf{Sí} &
  Solo infraestructura; desarrollo es one-time. \\ \rowcolor{altrow}
Toast POS + Mobile Order~\&~Pay & \$2,500 & \$90,000 & No &
  Sin mesa virtual con contrase\~na. \\
Square for Restaurants + KDS & \$2,200 & \$79,200 & No &
  Split bill desde cajero, no desde comensal. \\ \rowcolor{altrow}
Lightspeed Restaurant & \$3,500 & \$126,000 & No &
  Sin jerarqu\'ia cadena/zona/sucursal. \\
Yumminn & \$2,000 & \$72,000 & No &
  Sin KDS propio ni cuentas individuales. \\ \rowcolor{altrow}
Oracle MICROS POS & \$15,000 & \$540,000 & No &
  Enterprise. Hardware adicional requerido. \\ \rowcolor{green}
\textbf{Ahorro anual vs.\ promedio SaaS\,(excl.\ Oracle)} & & \textbf{\$81,360 MXN/a\~no} & & \\
\hline
\end{tabular}
\end{table}

\noindent El costo de infraestructura mensual del sistema (\$914 MXN total para todas
las sucursales) es significativamente inferior a cualquier alternativa SaaS del
mercado para una operaci\'on de 3~sucursales. El presupuesto total del proyecto
(\$1,198,241~MXN) se recupera contra las alternativas SaaS en aproximadamente
14~meses de operaci\'on, con la ventaja adicional de la propiedad total del c\'odigo
y sin costos recurrentes de licenciamiento.

\noindent\textbf{Veredicto:} El proyecto es econ\'omicamente viable. La inversi\'on
total del A\~no~1 (\$1,198,241~MXN / \$68,471~USD) es competitiva frente a las
alternativas SaaS del mercado y genera propiedad total del c\'odigo, ausencia de costos
recurrentes por licencias de plataforma y la posibilidad de escalar a otras cadenas
restauranteras. El detalle completo con f\'ormulas editables se encuentra en el archivo
\textit{Gesti\'on\_Costos.xlsx} adjunto. \textbf{FACTIBLE}.

% -----------------------------------------------------------
\FloatBarrier
\subsection{Factibilidad Legal}
% -----------------------------------------------------------

La factibilidad legal evalúa el cumplimiento del marco normativo mexicano e
internacional aplicable a la plataforma, incluyendo protección de datos personales,
pagos electrónicos, facturación electrónica, licencias de software y normativas
de telecomunicaciones.

\subsubsection{Ley Federal de Protección de Datos Personales en Posesión de
Particulares (LFPDPPP)}

La plataforma recolecta y procesa datos personales de los usuarios colaboradores
(correo electrónico, nombre, historial de accesos) a través de Supabase Auth.
Adicionalmente, se registran datos de comportamiento de comensales (órdenes
generadas, montos pagados) aunque de forma no nominal si el comensal no se
registra formalmente. Este tratamiento está sujeto a la LFPDPPP, publicada en
el Diario Oficial de la Federación (DOF) el 5 de julio de 2010, y a su Reglamento,
publicado el 21 de diciembre de 2011 \cite{covermanager_pricing, lfpdppp}.

Las sanciones por incumplimiento van desde 100 hasta 320,000 días de salario
mínimo general vigente (SMGV), equivalentes a aproximadamente \$12,200 hasta
\$39,040,000 MXN con el SMGV 2024 de \$248.93 MXN/día \cite{reglamento_lfpdppp}.

La Tabla~\ref{tab:lfpdppp} detalla las obligaciones aplicables y su estado de
cumplimiento en la versión 1.0 de la plataforma.

\setstretch{1.2}
\setlength{\tabcolsep}{4pt}
\begin{longtable}{>{\raggedright\arraybackslash}p{3.5cm}
                >{\centering\arraybackslash}p{1.5cm}
                >{\raggedright\arraybackslash}p{6.5cm}
                >{\raggedright\arraybackslash}p{4cm}}
\caption{Obligaciones LFPDPPP Aplicables a la Plataforma} \label{tab:lfpdppp} \\
\hline
\thead{Obligación legal} & \thead{Artículo} & \thead{Acción requerida} & \thead{Estado en v1.0} \\
\hline
\endfirsthead
\caption[]{TABLA \thetable{} (Continuación)} \\
\hline
\thead{Obligación legal} & \thead{Artículo} & \thead{Acción requerida} & \thead{Estado en v1.0} \\
\hline
\endhead
\hline
\endfoot
Aviso de Privacidad &
  Art.\ 15--16 &
  Redactar y publicar el Aviso de Privacidad antes del primer tratamiento de datos. Incluir identidad del responsable, finalidades, derechos ARCO y transferencias a Supabase/Vercel. &
  \textbf{Requerido} antes del lanzamiento. \\ \rowcolor{altrow}
Consentimiento del titular &
  Art.\ 8 &
  Obtener consentimiento expreso al crear cuenta o acceder por primera vez. &
  Implementar checkbox de aceptación en el registro. \\
Principio de licitud &
  Art.\ 7 &
  Los datos solo pueden tratarse para las finalidades declaradas en el Aviso de Privacidad. &
  Cumplido por diseño: datos usados solo para operar el servicio. \\ \rowcolor{altrow}
Principio de minimización &
  Art.\ 7 &
  Solo recolectar los datos estrictamente necesarios. &
  Cumplido: solo correo y nombre de empleados. Comensales son anónimos. \\
Derechos ARCO &
  Art.\ 28--37 &
  Implementar mecanismo para acceso, rectificación, cancelación y oposición. Plazo: 20 días hábiles. &
  Módulo de eliminación de cuenta requerido. \\ \rowcolor{altrow}
Medidas de seguridad &
  Art.\ 19 &
  TLS, cifrado en tránsito y en reposo, RLS en PostgreSQL, acceso mínimo por rol. &
  Implementado en la arquitectura del sistema. \\
Transferencias internacionales &
  Art.\ 36 &
  Datos alojados en Supabase (AWS us-east) están fuera de México. Debe declararse en el Aviso. &
  Requiere cláusula específica en el Aviso de Privacidad. \\ \rowcolor{altrow}
Responsable de datos &
  Art.\ 3 fracc.\ XIV &
  Designar un responsable de datos (Data Controller) dentro de la organización. &
  Designar antes del lanzamiento. \\
\end{longtable}
\setstretch{1.5}

\subsubsection{Estándar PCI DSS v4.0 y Decisión de Alcance de Pagos}

La plataforma \textbf{no integra en su versi\'on~1.0 un agregador de pago v\'ia tarjeta}.
Esta decisi\'on constituye una medida de gesti\'on de riesgo deliberada: procesar cobros con
tarjeta sin cumplir el est\'andar PCI~DSS (Payment Card Industry Data Security Standard)
versi\'on~4.0, publicado por el PCI Security Standards Council en marzo de~2022, expone a la
empresa a sanciones contractuales de las redes de pago (Visa, Mastercard) y a responsabilidad
civil ante fraudes \cite{conasami_2024}. El est\'andar es de cumplimiento obligatorio para cualquier entidad
que transmita, procese o almacene datos de titulares de tarjeta, y su incumplimiento puede
derivar en multas, revocaci\'on de la capacidad de aceptar tarjetas y da\~no reputacional.

La integraci\'on de cobros con tarjeta queda expl\'icitamente documentada como trabajo futuro
(LM-13). Cuando se incorpore en una versi\'on posterior, la estrategia recomendada ser\'a la
\textit{reducci\'on de alcance PCI} (scope reduction), delegando completamente el manejo de
datos de tarjeta a un proveedor certificado Level~1 PCI~DSS mediante tokenizaci\'on, de modo
que la plataforma nunca vea, transmita ni almacene datos de tarjeta.

La Tabla~\ref{tab:gateways} compara las opciones de gateways certificados disponibles en
M\'exico para cuando se eval\'ue esta integraci\'on en versiones futuras.

\begin{table}[H]
\caption{Gateways de Pago Certificados PCI DSS Level 1 en México}
\label{tab:gateways}
\centering
\small
\setlength{\tabcolsep}{3pt}
\begin{tabular}{>{\raggedright\arraybackslash}p{3cm}
                >{\raggedright\arraybackslash}p{3.5cm}
                >{\raggedright\arraybackslash}p{4.5cm}
                >{\raggedright\arraybackslash}p{5cm}}
\hline
\thead{Proveedor} & \thead{Certificación} & \thead{Comisión (México)} & \thead{Integración} \\
\hline
Stripe & Level 1 PCI DSS & 2.9\% + \$3~MXN (nacional) / 3.6\% + \$3~MXN (int.) & SDK oficial para Next.js (Stripe.js) \cite{pci_dss_4}. \\ \rowcolor{altrow}
Conekta & Level 1 PCI DSS & 2.8\% + \$3.5~MXN + IVA & SDK para Node.js \cite{stripe_pricing}. \\
OpenPay (BBVA) & Level 1 PCI DSS & 2.8\% + \$3~MXN + IVA & REST API + SDK JS. \\ \rowcolor{altrow}
Clip & Level 1 PCI DSS & 3.6\% (tarjeta no presente) + IVA & REST API. \\
\hline
\end{tabular}
\end{table}

\subsubsection{Obligaciones Fiscales: CFDI}

El Artículo 29 del Código Fiscal de la Federación (CFF) establece la obligación
de expedir Comprobantes Fiscales Digitales por Internet (CFDI) por los actos o
actividades de los contribuyentes \cite{conekta_pricing}. La plataforma en su versión 1.0 no emite
CFDI directamente; la obligación recae sobre la cadena restaurantera. No obstante,
se recomienda evaluar la integración con un Proveedor Autorizado de Certificación
(PAC) del SAT \cite{cff_art29} para automatizar la emisión de CFDI a solicitud del comensal,
especialmente en el segmento empresarial donde el 78\% de los comensales
de restaurantes de servicio completo requieren factura \cite{inegi_canirac_2023}.

\subsubsection{Ley Federal de Protección al Consumidor (LFPC)}

La plataforma intermedia en transacciones de consumo, lo que la sitúa dentro
del ámbito de la LFPC \cite{sat_pac}. Su Artículo 76 BIS regula el comercio electrónico,
estableciendo obligaciones de transparencia en precios, métodos de pago, términos
y condiciones, y el derecho a cancelar la operación antes de confirmar el pago.
La Tabla~\ref{tab:lfpc} detalla las obligaciones aplicables y su implementación
en la plataforma.

\begin{table}[H]
\caption{Obligaciones LFPC Aplicables a la Plataforma}
\label{tab:lfpc}
\centering
\small
\setlength{\tabcolsep}{3pt}
\begin{tabular}{>{\raggedright\arraybackslash}p{4cm}
                >{\centering\arraybackslash}p{2cm}
                >{\raggedright\arraybackslash}p{10cm}}
\hline
\thead{Obligación LFPC} & \thead{Artículo} & \thead{Implementación en la plataforma} \\
\hline
Información clara sobre precios antes de confirmar la compra &
  Art.\ 7 &
  La pantalla de pagos muestra el desglose completo del monto antes de solicitar datos de pago. \\ \rowcolor{altrow}
Confirmación explícita del pedido &
  Art.\ 76 BIS fracc.\ II &
  Diálogo de confirmación antes de enviar cada orden. Confirmación final antes de procesar el pago. \\
Derecho a cancelar antes de confirmar &
  Art.\ 76 BIS fracc.\ III &
  El comensal puede vaciar su canasta o cancelar el flujo de pago antes de confirmar la transacción. \\ \rowcolor{altrow}
Términos y condiciones accesibles &
  Art.\ 76 BIS fracc.\ I &
  Términos de Uso publicados y aceptados al acceder por primera vez a la plataforma. \\
Mecanismo de quejas &
  Art.\ 24 &
  Publicar correo o formulario de atención al cliente. PROFECO puede recibir quejas sobre el servicio. \\
\hline
\end{tabular}
\end{table}

\subsubsection{Licencias de Software del Stack}

El análisis de licencias garantiza que el uso comercial de cada componente del
stack sea legal. La Tabla~\ref{tab:licencias} muestra que ningún componente impone
restricciones incompatibles con el modelo de negocio SaaS de la plataforma \cite{lfpc, osi_licenses}.

\begin{table}[H]
\caption{Análisis de Licencias del Stack Tecnológico}
\label{tab:licencias}
\centering
\small
\setlength{\tabcolsep}{3pt}
\begin{tabular}{>{\raggedright\arraybackslash}p{3cm}
                >{\raggedright\arraybackslash}p{3cm}
                >{\centering\arraybackslash}p{2.5cm}
                >{\centering\arraybackslash}p{2.5cm}
                >{\raggedright\arraybackslash}p{5cm}}
\hline
\thead{Componente} & \thead{Licencia} & \thead{¿Uso comercial?} &
\thead{¿Publicar código?} & \thead{Restricciones relevantes} \\
\hline
Next.js & MIT License & Sí & No & Mantener aviso de copyright. \\ \rowcolor{altrow}
React & MIT License & Sí & No & Mantener aviso de copyright. \\
TypeScript & Apache 2.0 & Sí & No & Incluir copia de la licencia en distribuciones. \\ \rowcolor{altrow}
Prisma ORM & Apache 2.0 & Sí & No & Incluir copia de la licencia en distribuciones. \\
Supabase (cliente JS) & MIT License & Sí & No & Mantener aviso de copyright. \\ \rowcolor{altrow}
PostgreSQL & PostgreSQL License & Sí & No & Sin restricciones comerciales. \\
Node.js & MIT License & Sí & No & Mantener aviso de copyright. \\ \rowcolor{altrow}
Vercel (plataforma) & Propietario (SaaS) & Sí (Plan Pro) & N/A & Sujeto a Términos de Servicio de Vercel. \\
Supabase (plataforma) & Propietario + AGPLv3 & Sí (Plan Pro cloud) & Solo en self-hosting & En modo BaaS cloud, el código cliente no está sujeto a AGPL \cite{osi_licenses}. \\
\hline
\end{tabular}
\end{table}

\subsubsection{Normativas de Telecomunicaciones}

La plataforma es un servicio de software entregado a través de internet (SaaS/PWA)
y está regulada por la Ley Federal de Telecomunicaciones y Radiodifusión (LFTR) \cite{supabase_license}.
La plataforma no opera como concesionaria de red pública ni como proveedor de
servicios de telecomunicaciones (voz, SMS), por lo que la LFTR no impone
obligaciones de registro o concesión. No obstante, si en versiones futuras se
implementan notificaciones push vía SMS, se deberá revisar el Art. 191 de la
LFTR relativo a la identidad del remitente en comunicaciones comerciales.

\noindent\textbf{Veredicto:} El proyecto es legalmente viable. Las acciones legales
no negociables antes del lanzamiento son: (1) redacci\'on y publicaci\'on del Aviso de
Privacidad (Arts.~15--16 LFPDPPP); (2) implementaci\'on de consentimiento expreso y
mecanismo ARCO; y (3) T\'erminos y Condiciones conforme al Art.~76~BIS de la LFPC.
La decisi\'on de no integrar un agregador de pago v\'ia tarjeta en v1.0 elimina el riesgo
RL-02 (PCI~DSS) de la plataforma en esta versi\'on. El stack utiliza exclusivamente
licencias permisivas (MIT, Apache~2.0, BSD) sin restricciones sobre el uso comercial.
\textbf{FACTIBLE CON ACCIONES REQUERIDAS}.

% -----------------------------------------------------------
\subsection{Análisis y Mitigación de Riesgos}
% -----------------------------------------------------------

\subsubsection{Metodología de Evaluación}

Se identificaron y evaluaron 28 riesgos distribuidos en tres dimensiones complementarias:
\textbf{técnica} (seguridad de datos, integridad transaccional, concurrencia, disponibilidad
y calidad del software), \textbf{operativa} (errores humanos, capacitación del personal y
continuidad del negocio) y \textbf{legal} (LFPDPPP, CFF/SAT, PCI DSS, LFPC y
responsabilidad civil) \cite{covermanager_pricing, lfpdppp, conasami_2024, sat_pac}.

Cada riesgo fue calificado conforme a la escala que se describe a continuación. La
probabilidad refleja la frecuencia esperada de ocurrencia: \textit{Alta} ($>$50\%),
\textit{Media} (20--50\%) o \textit{Baja} ($<$20\%). El impacto pondera las consecuencias
sobre la operaci\'on, los datos personales y la exposici\'on legal: \textit{Alto}, \textit{Medio}
o \textit{Bajo}. El nivel de riesgo compuesto se determina mediante la
Tabla~\ref{tab:matriz_compuesta}.

\begin{table}[H]
\caption{Matriz de Calificación Compuesta Probabilidad $\times$ Impacto}
\label{tab:matriz_compuesta}
\centering
\small
\setlength{\tabcolsep}{6pt}
\begin{tabular}{>{\bfseries\raggedright\arraybackslash}p{2.5cm}
                >{\centering\arraybackslash}p{2.5cm}
                >{\centering\arraybackslash}p{2.5cm}
                >{\centering\arraybackslash}p{2.5cm}}
\hline
\thead{Prob. \textbackslash\ Impacto} & \thead{Bajo} & \thead{Medio} & \thead{Alto} \\
\hline
Alta  & \cellcolor{orange}ALTO    & \cellcolor{orange}ALTO     & \cellcolor{red}CR\'ITICO \\
Media & \cellcolor{altrow}MEDIO   & \cellcolor{orange}ALTO     & \cellcolor{orange}ALTO \\
Baja  & \cellcolor{green}BAJO     & \cellcolor{altrow}MEDIO    & \cellcolor{orange}ALTO \\
\hline
\end{tabular}
\end{table}

Los tipos de tratamiento aplicados son: \textbf{Mitigar} (controles t\'ecnicos u operativos
para reducir probabilidad o impacto), \textbf{Transferir} (delegar el riesgo a un tercero),
\textbf{Aceptar} (documentar y monitorear cuando el costo de mitigaci\'on supera el impacto)
y \textbf{Evitar} (eliminar la actividad que genera el riesgo).

\subsubsection{Resumen Ejecutivo de Riesgos}

La distribuci\'on por nivel de riesgo se resume en la Tabla~\ref{tab:resumen_riesgos}.
Se identificaron 5~riesgos cr\'iticos que requieren resoluci\'on obligatoria antes del
lanzamiento de la plataforma.

\begin{table}[H]
\caption{Distribución de Riesgos por Nivel}
\label{tab:resumen_riesgos}
\centering
\small
\setlength{\tabcolsep}{5pt}
\begin{tabular}{>{\bfseries\centering\arraybackslash}p{2.5cm}
                >{\centering\arraybackslash}p{1.5cm}
                >{\raggedright\arraybackslash}p{5cm}
                >{\centering\arraybackslash}p{2.5cm}}
\hline
\thead{Nivel} & \thead{\# Riesgos} & \thead{Criterio de acción} & \thead{Prioridad} \\
\hline
\rowcolor{red}
CR\'ITICO & 5 & Acci\'on inmediata antes del lanzamiento. & Inmediata \\ \rowcolor{orange}
ALTO     & 17 & Atender en el sprint actual o pr\'oximo. & Alta \\ \rowcolor{altrow}
MEDIO    & 5  & Planificar mitigaci\'on en los primeros 2 meses post-lanzamiento. & Media \\ \rowcolor{green}
BAJO     & 1  & Monitorear; atender si la frecuencia aumenta. & Baja \\
\hline
\end{tabular}
\end{table}

\noindent Los hallazgos cr\'iticos que requieren atenci\'on antes del lanzamiento son:
RT-01/RT-02 (concurrencia en pagos sin control definido), RT-03 (transmisi\'on de credenciales
sin cifrado), RL-01 (recoleci\'on de datos sin aviso de privacidad, sancionable bajo LFPDPPP)
y RO-01 (inexistencia de flujo de recuperaci\'on de contrase\~na de mesa).

\noindent\textit{Nota importante --- alcance de pagos en v1.0:} La versi\'on~1.0 de la
plataforma \textbf{no contempla la integraci\'on de un agregador de pago v\'ia tarjeta}.
Esta decisi\'on es una medida de gesti\'on de riesgo deliberada: procesar pagos con tarjeta
sin cumplir PCI~DSS~v4.0 expone a la empresa a sanciones contractuales de las redes de
pago (Visa, Mastercard) y a responsabilidad civil ante fraudes \cite{conasami_2024}. El riesgo RL-02 queda
por tanto \textbf{evitado} en esta versi\'on; su integraci\'on en versiones futuras deber\'a
ir precedida de la certificaci\'on correspondiente o de la contrataci\'on de un proveedor
Level~1 PCI~DSS que asuma el alcance del entorno de datos de tarjeta (CDE).

\subsubsection{Riesgos Técnicos (RT)}

La Tabla~\ref{tab:riesgos_tecnicos} presenta los 11~riesgos t\'ecnicos identificados,
ordenados por nivel de criticidad.

\setstretch{1.2}
\setlength{\tabcolsep}{3pt}
\begin{longtable}{>{\bfseries\centering\arraybackslash}p{0.7cm}
                >{\raggedright\arraybackslash}p{4cm}
                >{\centering\arraybackslash}p{1.1cm}
                >{\centering\arraybackslash}p{1.1cm}
                >{\centering\arraybackslash}p{1.4cm}
                >{\raggedright\arraybackslash}p{6cm}
                >{\centering\arraybackslash}p{1.2cm}}
\caption{Matriz de Riesgos Técnicos} \label{tab:riesgos_tecnicos} \\
\hline
\thead{ID} & \thead{Descripción del riesgo} & \thead{Prob.} &
\thead{Impacto} & \thead{Nivel} & \thead{Estrategia de mitigación} & \thead{Trat.} \\
\hline
\endfirsthead
\caption[]{TABLA \thetable{} (Continuación)} \\
\hline
\thead{ID} & \thead{Descripción del riesgo} & \thead{Prob.} &
\thead{Impacto} & \thead{Nivel} & \thead{Estrategia de mitigación} & \thead{Trat.} \\
\hline
\endhead
\hline
\endfoot
\rowcolor{red}
RT-01 &
  Concurrencia en liquidaci\'on: dos comensales pagan la cuenta del mismo comensal
  simult\'aneamente. &
  Alta & Alto & CR\'ITICO &
  Bloqueo optimista (versioning) o pesimista (mutex por ID de comensal).
  Validar saldo pendiente al aplicar el pago, no al iniciar el flujo. &
  Mitigar \\ \rowcolor{red}
RT-02 &
  Sobrepago en mesa compartida: aportes simult\'aneos superan el monto total pendiente. &
  Alta & Alto & CR\'ITICO &
  Transacciones at\'omicas con bloqueo a nivel de fila (\texttt{SELECT FOR UPDATE}).
  Retornar error al segundo pago si el monto ya fue cubierto.
  Implementar \texttt{idempotency\_keys} por operaci\'on. &
  Mitigar \\ \rowcolor{orange}
RT-03 &
  Brecha de seguridad: credenciales de mesa virtual interceptadas en tr\'ansito
  (ID + contrase\~na). &
  Media & Alto & ALTO &
  Forzar HTTPS/TLS~1.2+ en todas las comunicaciones. Almacenar contrase\~na
  hasheada (bcrypt). Limitar la vida \'util de la contrase\~na al ciclo de la mesa. &
  Mitigar \\ \rowcolor{orange}
RT-04 &
  Fuerza bruta sobre contrase\~na de mesa virtual: intentos masivos de acceso. &
  Media & Medio & ALTO &
  Rate limiting por IP y por ID de mesa (m\'ax.\ 5 intentos antes de bloqueo
  temporal). Registrar intentos fallidos para audit\'oria \cite{mdn_websocket}. &
  Mitigar \\ \rowcolor{orange}
RT-05 &
  Inconsistencia de datos de orden: el Chef visualiza una orden cancelada sin
  actualizaci\'on en tiempo real. &
  Media & Alto & ALTO &
  WebSockets v\'ia Supabase~Realtime para el tablero del Chef \cite{supabase_license}.
  Estados de orden con transiciones at\'omicas e irreversibles. &
  Mitigar \\ \rowcolor{orange}
RT-06 &
  Degradaci\'on del flujo de liquidaci\'on manual: el comensal no puede completar
  su registro de pago por fallo de conectividad. &
  Baja & Alto & ALTO &
  Implementar reintentos autom\'aticos con backoff exponencial.
  Mostrar mensaje de error con instrucci\'on de pago alternativo en caja. &
  Mitigar \\ \rowcolor{orange}
RT-07 &
  Gesti\'on de sesiones deficiente: sesi\'on del anfitri\'on expira bloqueando la
  mesa. &
  Media & Medio & ALTO &
  Refresh tokens con vida prolongada (m\'inimo igual al turno del restaurante).
  Permitir reingreso por ID + contrase\~na sin re-escaneo de QR. &
  Mitigar \\ \rowcolor{orange}
RT-08 &
  P\'erdida o corrupci\'on de datos transaccionales por ca\'ida de base de datos
  durante liquidaci\'on activa. &
  Baja & Alto & ALTO &
  WAL (Write-Ahead Logging) y r\'eplicas en PostgreSQL/Supabase. Transacciones
  ACID. Respaldos autom\'aticos. RTO $<$ 4~hrs \cite{mdn_pwa}. &
  Mitigar \\ \rowcolor{altrow}
RT-09 &
  C\'odigo QR clonado o reutilizado: un QR de sesi\'on previa permite acceso
  no autorizado. &
  Baja & Medio & MEDIO &
  QR con firma HMAC-SHA256 v\'alida que impida la manipulaci\'on de enlaces.
  Invalidar el token en el backend al activarse. &
  Mitigar \\ \rowcolor{orange}
RT-10 &
  Escalabilidad insuficiente en horas pico: degradaci\'on con m\'ultiples mesas
  activas. &
  Media & Medio & ALTO &
  Pruebas de carga antes del despliegue (k6, JMeter). Cach\'e en lecturas del
  men\'u. Arquitectura serverless con escalado horizontal (Vercel + Supabase) \cite{vercel_pricing, supabase_pricing}. &
  Mitigar \\ \rowcolor{orange}
RT-11 &
  Inyecci\'on de datos maliciosos: cantidades negativas o caracteres especiales en
  campos de orden. &
  Media & Medio & ALTO &
  Validar y sanitizar todos los inputs en el backend. Consultas parametrizadas
  (Prisma ORM previene SQL Injection por dise\~no). L\'imites de cantidad (m\'in.\ 1,
  m\'ax.\ configurable) \cite{prisma_docs}. &
  Mitigar \\
\end{longtable}
\setstretch{1.5}

\subsubsection{Riesgos Operativos (RO)}

La Tabla~\ref{tab:riesgos_operativos} presenta los 9~riesgos operativos, derivados de
procesos, comportamiento humano y procedimientos de operaci\'on diaria.

\setstretch{1.2}
\setlength{\tabcolsep}{3pt}
\begin{longtable}{>{\bfseries\centering\arraybackslash}p{0.7cm}
                >{\raggedright\arraybackslash}p{4cm}
                >{\centering\arraybackslash}p{1.1cm}
                >{\centering\arraybackslash}p{1.1cm}
                >{\centering\arraybackslash}p{1.4cm}
                >{\raggedright\arraybackslash}p{6cm}
                >{\centering\arraybackslash}p{1.2cm}}
\caption{Matriz de Riesgos Operativos} \label{tab:riesgos_operativos} \\
\hline
\thead{ID} & \thead{Descripción del riesgo} & \thead{Prob.} &
\thead{Impacto} & \thead{Nivel} & \thead{Estrategia de mitigación} & \thead{Trat.} \\
\hline
\endfirsthead
\caption[]{TABLA \thetable{} (Continuación)} \\
\hline
\thead{ID} & \thead{Descripción del riesgo} & \thead{Prob.} &
\thead{Impacto} & \thead{Nivel} & \thead{Estrategia de mitigación} & \thead{Trat.} \\
\hline
\endhead
\hline
\endfoot
\rowcolor{red}
RO-01 &
  Anfitri\'on olvida contrase\~na de mesa y ning\'un comensal la conoce: mesa queda
  inaccesible con \'ordenes activas. &
  Alta & Alto & CR\'ITICO &
  Protocolo operativo: el mesero puede visualizar el identificador y generar una
  contrase\~na de emergencia desde su rol. Flujo de recuperaci\'on incluido en el
  roadmap (v1.1). &
  Aceptar \\ \rowcolor{orange}
RO-02 &
  Mesero inicializa una mesa virtual con mesas f\'isicas incorrectas (asociaci\'on
  err\'onea). &
  Media & Medio & ALTO &
  Confirmaci\'on visual con n\'umero de mesa antes de confirmar la inicializaci\'on.
  Permitir correcci\'on antes de que el anfitri\'on escanee el QR.
  Registrar log de inicializaci\'on. &
  Mitigar \\ \rowcolor{orange}
RO-03 &
  Personal sin capacitaci\'on adecuada: errores operativos desde el primer d\'ia. &
  Alta & Medio & ALTO &
  Manual de usuario por rol. Onboarding interactivo en la app (tooltips, tutorial
  guiado). Modo sandbox para pr\'actica. Superusuario de sucursal para soporte
  de primer nivel. &
  Mitigar \\ \rowcolor{orange}
RO-04 &
  Men\'u desactualizado: comensales ordenan productos descontinuados o con precios
  incorrectos. &
  Media & Medio & ALTO &
  Sincronizaci\'on peri\'odica entre el men\'u f\'isico y la plataforma. Acceso del
  gerente para actualizar en tiempo real. Alertas si un producto no ha sido
  actualizado en m\'as de 30~d\'ias. &
  Mitigar \\ \rowcolor{orange}
RO-05 &
  Chef no recibe notificaci\'on de nuevas \'ordenes en tiempo real, generando
  retrasos en cocina. &
  Media & Alto & ALTO &
  Notificaciones push o sonoras ante nuevas \'ordenes (Web Audio API).
  SLA de respuesta: confirmaci\'on de orden en $<$60~s. Alerta al gerente si una
  orden no es confirmada a tiempo \cite{supabase_license}. &
  Mitigar \\ \rowcolor{altrow}
RO-06 &
  Mesa compartida abandonada: comensales se van sin liquidar su parte. &
  Baja & Medio & MEDIO &
  Mesero y gerente visualizan el estatus de liquidaci\'on en tiempo real.
  Protocolo de cobro en caja como respaldo. Decisi\'on de negocio sobre bloqueo
  de mesa hasta liquidaci\'on total. &
  Aceptar \\ \rowcolor{altrow}
RO-07 &
  Dos personas intentan activar la misma mesa virtual como anfitri\'on
  simult\'aneamente. &
  Media & Bajo & MEDIO &
  Activaci\'on de mesa permitida una sola vez (PENDIENTE~$\rightarrow$~ACTIVA).
  El segundo intento retorna error indicando que la mesa ya fue activada. &
  Mitigar \\ \rowcolor{altrow}
RO-08 &
  Credenciales de mesa compartidas con personas externas al grupo. &
  Baja & Medio & MEDIO &
  Advertencia expl\'icita al anfitri\'on sobre no compartir la contrase\~na.
  L\'imite m\'aximo configurable de comensales por mesa.
  Mesero puede auditar el n\'umero de comensales activos. &
  Aceptar \\ \rowcolor{green}
RO-09 &
  El mesero olvida cambiar el estatus de la mesa de Liquidada a Libre. &
  Baja & Bajo & BAJO &
  Notificaci\'on autom\'atica al mesero/gerente cuando la mesa pasa a estado
  Liquidada. Transici\'on autom\'atica a Libre tras tiempo configurable de
  inactividad post-liquidaci\'on. &
  Mitigar \\
\end{longtable}
\setstretch{1.5}

\subsubsection{Riesgos Legales y Regulatorios (RL)}

El marco normativo aplicable comprende la LFPDPPP \cite{covermanager_pricing, lfpdppp}, el C\'odigo Fiscal de la
Federaci\'on (CFF) \cite{conekta_pricing}, el est\'andar PCI~DSS~v4.0 \cite{conasami_2024}, la Ley Federal de Protecci\'on al
Consumidor (LFPC) \cite{sat_pac} y el C\'odigo Civil Federal. La Tabla~\ref{tab:riesgos_legales}
presenta los 8~riesgos legales identificados.

\setstretch{1.2}
\setlength{\tabcolsep}{3pt}
\begin{longtable}{>{\bfseries\centering\arraybackslash}p{0.7cm}
                >{\raggedright\arraybackslash}p{4cm}
                >{\centering\arraybackslash}p{1.1cm}
                >{\centering\arraybackslash}p{1.1cm}
                >{\centering\arraybackslash}p{1.4cm}
                >{\raggedright\arraybackslash}p{6cm}
                >{\centering\arraybackslash}p{1.2cm}}
\caption{Matriz de Riesgos Legales y Regulatorios} \label{tab:riesgos_legales} \\
\hline
\thead{ID} & \thead{Descripción del riesgo} & \thead{Prob.} &
\thead{Impacto} & \thead{Nivel} & \thead{Estrategia de mitigación} & \thead{Trat.} \\
\hline
\endfirsthead
\caption[]{TABLA \thetable{} (Continuación)} \\
\hline
\thead{ID} & \thead{Descripción del riesgo} & \thead{Prob.} &
\thead{Impacto} & \thead{Nivel} & \thead{Estrategia de mitigación} & \thead{Trat.} \\
\hline
\endhead
\hline
\endfoot
\rowcolor{red}
RL-01 &
  Incumplimiento LFPDPPP: recoleci\'on de datos personales sin aviso de privacidad
  ni consentimiento expl\'icito. Sanciones de hasta 320,000 d\'ias de SMGV \cite{reglamento_lfpdppp}. &
  Alta & Alto & CR\'ITICO &
  Redactar y publicar Aviso de Privacidad (Art.~16 LFPDPPP). Consentimiento expreso
  al registrarse. Nombrar un Data Controller. Definir plazos de retenci\'on y
  procedimiento de eliminaci\'on (Art.~19) \cite{covermanager_pricing}. &
  Mitigar \\ \rowcolor{orange}
RL-02 &
  PCI~DSS: procesamiento de pagos con tarjeta sin certificaci\'on adecuada, con
  exposici\'on a sanciones de redes de pago y responsabilidad civil ante fraudes \cite{conasami_2024}. &
  --- & Alto & EVITADO &
  \textbf{Decisión de alcance v1.0:} la plataforma \textbf{no integra} agregador de
  pago v\'ia tarjeta. Este riesgo queda \textbf{eliminado} en la versi\'on actual.
  La integraci\'on futura requerir\'a proveedor Level~1 PCI~DSS o auto-certificaci\'on. &
  Evitar \\ \rowcolor{orange}
RL-03 &
  Obligaci\'on fiscal SAT: transacciones de pago podr\'ian requerir emisi\'on de CFDI
  conforme al CFF Art.~29 \cite{conekta_pricing}. &
  Media & Alto & ALTO &
  Consultor\'ia fiscal para determinar responsabilidad de emisi\'on de CFDI
  (cadena restaurantera vs.\ plataforma). Evaluar integraci\'on con PAC
  autorizado por el SAT \cite{cff_art29} en versiones futuras. &
  Mitigar \\ \rowcolor{orange}
RL-04 &
  Responsabilidad por cobros no autorizados o err\'oneos: disputa del comensal
  ante PROFECO o banco (chargeback). &
  Media & Alto & ALTO &
  Comprobante de pago digital inmediato. Registro detallado e inmutable de
  transacciones (qu\'e, a qui\'en, monto, timestamp). Pol\'itica de devoluciones
  publicada en Términos de Servicio \cite{sat_pac}. &
  Mitigar \\ \rowcolor{orange}
RL-05 &
  Ausencia de T\'erminos y Condiciones de Uso: plataforma opera sin contratos
  que delimiten responsabilidades entre partes. &
  Media & Medio & ALTO &
  Redactar T\'erminos y Condiciones aceptados al acceder por primera vez.
  Especificar limitaci\'on de responsabilidad, jurisdicci\'on (M\'exico) y
  mecanismo de quejas \cite{sat_pac}. &
  Mitigar \\ \rowcolor{orange}
RL-06 &
  Acceso de menores de edad a la plataforma sin consentimiento de padres o tutores. &
  Baja & Alto & ALTO &
  T\'erminos de Uso restringen el acceso a mayores de 18~a\~nos. Confirmaci\'on de
  mayor\'ia de edad al registrarse. Supervisi\'on operativa a trav\'es del
  anfitri\'on adulto. &
  Mitigar \\ \rowcolor{orange}
RL-07 &
  Violaci\'on de derechos ARCO: usuario solicita eliminar sus datos y la
  plataforma no tiene mecanismo (Arts.~28--37 LFPDPPP) \cite{covermanager_pricing}. &
  Baja & Alto & ALTO &
  Funcionalidad de eliminaci\'on de cuenta y datos asociados.
  Procedimiento ARCO publicado en la app.
  Respuesta en plazo legal de 20~d\'ias h\'abiles. &
  Mitigar \\ \rowcolor{altrow}
RL-08 &
  Propiedad intelectual: uso de librer\'ias con licencias incompatibles con uso
  comercial (GPL, AGPL). &
  Baja & Medio & MEDIO &
  Auditor\'ia de licencias del stack antes del lanzamiento (FOSSA, licensecheck).
  Preferir MIT, Apache~2.0, BSD. Documentar todas las dependencias \cite{lfpc, osi_licenses}. &
  Mitigar \\
\end{longtable}
\setstretch{1.5}

\subsubsection{Hoja de Ruta de Mitigación}

La Tabla~\ref{tab:hoja_ruta} define el orden de atenci\'on recomendado para los riesgos
identificados, priorizando los de nivel Cr\'itico y Alto.

\setstretch{1.2}
\setlength{\tabcolsep}{4pt}
\begin{longtable}{>{\bfseries\raggedright\arraybackslash}p{3cm}
                >{\raggedright\arraybackslash}p{3cm}
                >{\raggedright\arraybackslash}p{9.5cm}}
\caption{Hoja de Ruta de Mitigación de Riesgos} \label{tab:hoja_ruta} \\
\hline
\thead{Etapa} & \thead{Riesgos a atender} & \thead{Acciones principales} \\
\hline
\endfirsthead
\caption[]{TABLA \thetable{} (Continuación)} \\
\hline
\thead{Etapa} & \thead{Riesgos a atender} & \thead{Acciones principales} \\
\hline
\endhead
\hline
\endfoot
Antes del lanzamiento &
  RT-01, RT-02, RT-03, RT-09, RL-01, RO-01 &
  Control de concurrencia (\texttt{SELECT FOR UPDATE}), TLS/HTTPS, QR de un
  solo uso con firma HMAC, Aviso de Privacidad (LFPDPPP), protocolo de emergencia
  para mesa sin contrase\~na. \\ \rowcolor{altrow}
Sprint~1 post-lanzamiento (sem.\ 1--4) &
  RT-04, RT-05, RT-07, RT-11, RL-03, RL-04, RL-05 &
  Rate limiting, Supabase Realtime para tablero del Chef, gesti\'on de sesiones,
  validaci\'on de inputs, consultor\'ia fiscal SAT/CFDI, comprobante de pago
  digital, T\'erminos y Condiciones. \\
Sprint~2 post-lanzamiento (sem.\ 5--8) &
  RT-06, RT-10, RO-02, RO-03, RO-05, RL-06, RL-07 &
  Reintentos con backoff, pruebas de carga (k6), confirmaci\'on visual de mesas,
  capacitaci\'on y onboarding, notificaciones al Chef, restricci\'on de menores,
  mecanismo ARCO. \\ \rowcolor{altrow}
Mediano plazo (mes~3--6) &
  RT-08, RO-04, RO-06, RO-07, RO-08, RO-09, RL-08 &
  Plan de recuperaci\'on ante desastres (RTO $<$ 4~hrs), gesti\'on de men\'u en
  tiempo real, pol\'iticas de retenci\'on de datos, automatizaci\'on de
  transici\'on de estatus de mesa, auditor\'ia de licencias. \\
\end{longtable}
\setstretch{1.5}

\subsubsection{Conclusiones del Análisis de Riesgos}

\noindent\textbf{Factibilidad T\'ecnica.}
La plataforma es t\'ecnicamente factible. Los riesgos t\'ecnicos identificados son manejables
mediante patrones de ingenier\'ia est\'andar (control de concurrencia, cifrado, validaci\'on
de inputs, WebSockets). El principal riesgo de implementaci\'on radica en la gesti\'on de
transacciones concurrentes (RT-01, RT-02), que debe resolverse en la capa de backend antes
de entrar a producci\'on mediante \texttt{SELECT FOR UPDATE} en PostgreSQL.

\noindent\textbf{Factibilidad Operativa.}
La plataforma es operativamente viable siempre que se acompa\~ne de un programa de
capacitaci\'on por rol y se definan protocolos claros para escenarios de falla. El riesgo
operativo m\'as urgente (RO-01) requiere una decisi\'on de producto antes del lanzamiento:
implementar el flujo de recuperaci\'on de contrase\~na en v1.0, o documentar y capacitar
un protocolo manual que el mesero pueda ejecutar.

\noindent\textbf{Factibilidad Legal.}
El cumplimiento de la LFPDPPP es no negociable y debe resolverse antes de procesar
cualquier dato personal. El riesgo RL-02 (PCI~DSS) ha sido \textbf{eliminado} en v1.0
mediante la decisi\'on de no integrar un agregador de pago v\'ia tarjeta: procesar cobros
con tarjeta sin certificaci\'on adecuada expone a la empresa a sanciones contractuales de
las redes de pago y a responsabilidad civil ante fraudes, un riesgo que supera los
beneficios de incluir dicha funcionalidad en esta versi\'on \cite{conasami_2024}.

\noindent\textbf{Veredicto.}
El proyecto es factible en sus tres dimensiones, condicionado a la resoluci\'on de los
5~riesgos cr\'iticos antes del lanzamiento. Con las mitigaciones propuestas implementadas,
el perfil de riesgo residual de la plataforma es aceptable para una operaci\'on en producci\'on.

% ============================================================
%  CAPÍTULO 4 — METODOLOGÍA Y ANÁLISIS
% ============================================================
\newpage
